\begin{textsample}{POS Dim 2 – am – Score 26.00 – am/am\_em\_29.txt}  \label{ex:f2_pos_151}
3.4 Especificidades do Itinerário \textbf{Formativo} de Formação \textbf{Técnica} e Profissional As \textbf{Diretrizes} Curriculares Nacionais Gerais para a Educação Profissional Tecnológica ( DCNGEPT ), marco normativo definido pela Resolução CNE/CP nº 01/2021 e que traz os princípios e critérios a serem observados pelos sistemas e instituições de ensino públicas e privadas para o planejamento, organização, desenvolvimento, avaliação da EPT, presencial e a distância, especificamente, dispõe, em seu § 5º do art. 5º, que o Itinerário \textbf{Formativo} na Educação \textbf{Técnica} e Profissional é o conjunto de unidades curriculares, etapas ou módulos que compõe a sua organização em eixos tecnológicos e respectiva área tecnológica.
Desse modo, ao fazer o cotejo entre os dois marcos legais ( DCN e DCNGEPT ), o Itinerário \textbf{Formativo} na Educação \textbf{Técnica} e Profissional possibilita aprofundar conhecimentos e desenvolver habilidades associadas aos eixos \textbf{estruturantes}, às habilidades básicas requeridas pelo mundo do trabalho e às habilidades específicas relacionadas aos \textbf{Cursos} \textbf{Técnicos}, \textbf{Cursos} de \textbf{Qualificação} \textbf{Profissional} ( FICs ) ou Programa de Aprendizagem Profissional escolhidos pelos estudantes, a fim de prepará -los para o \textbf{prosseguimento} de estudos e para o mundo do trabalho de modo a contribuir, para a construção de soluções de problemas específicos da sociedade, numa perspectiva de formação humanística, emancipatória e integral do indivíduo.
Para a organização curricular e planejamento da \textbf{oferta} de Itinerário de Formação \textbf{Técnica} e Profissional articulada ao Ensino Médio é fundamental pontuar a necessidade de um diálogo constante com os marcos normativos que constituem a base para o planejamento e a \textbf{oferta} de \textbf{cursos} e programa de Educação Profissional e Tecnológica, como a Resolução CNE/CP nº 1/2021, o \textbf{Catálogo} Nacional de \textbf{Cursos} \textbf{Técnicos} ( CNCT ) e a Classificação Brasileira de Ocupações ( CBO ).
Para \textbf{oferta} de aprofundamento de \textbf{Itinerários} \textbf{Formativos} de Formação \textbf{Técnica} e Profissional, deverão ser observados os \textbf{requisitos} mínimos constantes na Resolução CNE/CP nº 1/2021.
Na estrutura do Novo Ensino Médio, os tipos de organização curricular dos \textbf{Itinerários} \textbf{Formativos} de Formação \textbf{Técnica} e Profissional poderão ser divididos nos seguintes \textbf{cursos} e programas : - \textbf{Curso} de Formação Inicial ou \textbf{Qualificação} \textbf{Profissional} ; - \textbf{Curso} \textbf{Técnico} de Nível Médio ; - Programa de Aprendizagem Profissional.
Na \textbf{qualificação} \textbf{profissional} ou formação inicial, os \textbf{cursos} e programas de educação \textbf{profissional} devem ter \textbf{carga} \textbf{horária} mínima de 160 horas e serão organizados em \textbf{trajetórias} de formação, que favoreçam o aproveitamento contínuo e articulado de estudos, conforme dispõe o Decreto nº 8.268/2014, que alterou o Decreto nº 5.154/2004.
Na Educação \textbf{Técnica} e \textbf{Profissional} de Nível Médio, as \textbf{ofertas} devem contemplar os eixos tecnológicos da Formação \textbf{Técnica} e Profissional, bem como observar a Classificação Brasileira de Ocupações ( CBO ), o Código e Título da Família Ocupacional, as normas associadas ao exercício \textbf{profissional} e as \textbf{cargas} \textbf{horárias} mínimas de 800, 1.000 e 1.200 horas, conforme \textbf{habilitação} \textbf{profissional} de \textbf{técnico}, conforme preveem os arts. 29 e 30 da Resolução CNE/CP nº 1/2021.
Ao serem organizados por eixos tecnológicos, os \textbf{cursos} e programas de Educação \textbf{Técnica} e Profissional possibilitam a construção de diferentes \textbf{itinerários} \textbf{formativos} e, conforme a LDB, em seu art. 36, § 6º, com a redação dada pela Lei nº 13.415/2017, deverão considerar também : I - a inclusão de vivências práticas de trabalho no setor produtivo ou em ambientes de simulação, estabelecendo parcerias e fazendo uso, quando aplicável, de instrumentos estabelecidos pela \textbf{legislação} sobre aprendizagem \textbf{profissional} ; II - a possibilidade de concessão de certificados intermediários de \textbf{qualificação} para o trabalho, quando a formação for estruturada e organizada em etapas com terminalidade ( BRASIL, 1996 ).
É importante destacar que, quando o \textbf{curso} ou programa de Educação \textbf{Técnica} e \textbf{Profissional} de Nível Médio for organizado por saídas intermediárias, correspondentes a etapas de \textbf{qualificação} \textbf{profissional} \textbf{técnica}, conforme prevê o art. 26 da Resolução CNE/CP nº 1/2021, a \textbf{carga} \textbf{horária} mínima, para cada etapa com terminalidade de \textbf{qualificação} \textbf{profissional} \textbf{técnica} prevista em um \textbf{itinerário} \textbf{formativo} de \textbf{curso} \textbf{técnico} é de 20 \% ( vinte por cento ) da \textbf{carga} \textbf{horária} mínima prevista para a respectiva \textbf{habilitação} \textbf{profissional}, indicada no \textbf{Catálogo} Nacional de \textbf{Cursos} \textbf{Técnicos}.
Quanto à forma da \textbf{oferta}, segundo a Resolução CNE/CP nº 1/2021, os \textbf{cursos} \textbf{técnicos} serão desenvolvidos nas formas : I - integrada, \textbf{ofertada} somente a quem já tenha concluído o Ensino Fundamental, com matrícula única na mesma instituição, de modo a conduzir o estudante à \textbf{habilitação} \textbf{profissional} \textbf{técnica} ao mesmo tempo em que conclui a última etapa da Educação Básica ; II - concomitante, \textbf{ofertada} a quem ingressa no Ensino Médio ou já o esteja \textbf{cursando}, \textbf{efetuando} -se matrículas distintas para cada \textbf{curso}, aproveitando oportunidades educacionais disponíveis, seja em unidades de ensino da mesma instituição ou em distintas instituições e \textbf{redes} de ensino ; III - concomitante intercomplementar, desenvolvida simultaneamente em distintas instituições ou \textbf{rede} de ensino, mais integrada no conteúdo, mediante a ação de convênio ou acordo de intercomplementaridade, para execução de projeto pedagógico unificado ; e IV - subsequente, desenvolvida em \textbf{cursos} \textbf{destinados} exclusivamente a quem já tenha concluído o Ensino Médio ( BRASIL, 2021 ).
Na organização do \textbf{itinerário} de \textbf{qualificação} \textbf{profissional} ( e também na \textbf{habilitação} \textbf{profissional} \textbf{técnica} ), podem ser incluídos os programas de aprendizagem \textbf{profissional}, nos termos da Lei nº 10.097/2000 ( Lei da Aprendizagem ).
O Sistema de Aprendizagem Profissional ou Programa de Aprendizagem Profissional é um percurso importante que une a ampliação da experiência \textbf{profissional} prática com a formação teórica para \textbf{adolescentes} e jovens de 14 a 24 anos que frequentem a Educação Básica, permitindo a articulação entre estudo e trabalho.
Como se pode perceber, o Programa de Aprendizagem Profissional é um importante tipo de \textbf{oferta} de Itinerário de Formação \textbf{Técnica} e Profissional e que pode ser articulado ao novo Ensino Médio, favorecendo o estudante que precisa trabalhar e também estudar.
Vale pontuar que os \textbf{cursos} devem seguir as nomenclaturas dispostas no \textbf{Catálogo} Nacional de \textbf{Cursos} \textbf{Técnicos} ( CNCT ), bem como trazer o código da Classificação Brasileira de Ocupações ( CBO ) associado ao \textbf{curso}.
Orienta -se, ainda, que a duração do programa seja de até dois anos, a começar, no mínimo, a partir do 2º \textbf{semestre} do \textbf{Curso}, com atividades práticas de 4 horas ou 6 horas por dia, até 5 dias por semana, limitada a 25 horas \textbf{semanais}.
Finalmente, destaca -se que, conforme Resolução CNE/CP nº 1/2021, as instituições e \textbf{redes} que oferecem Educação Profissional e Tecnológica poderão \textbf{ofertar} \textbf{cursos} experimentais que não constem no CNCT desde que devidamente autorizados pelos \textbf{órgãos} próprios dos respectivos sistemas de ensino e informem esta condição de \textbf{cursos} experimentais aos candidatos a esses \textbf{cursos}.
Na sequência, apresenta -se recomendações quanto à elaboração de \textbf{Itinerários} \textbf{Formativos}.

% matched lemmas: adolescente, carga, catálogo, cursar, curso, destinar, diretriz, efetuar, estruturante, formativo, habilitação, horário, itinerário, legislação, oferta, ofertar, profissional, prosseguimento, qualificação, rede, requisito, semanal, semestre, trajetória, técnico, órgão
\end{textsample}
